% Automatic exercise-number references: see \Solution in layout.tex.
% Each solution is indented as one coherent, multi-line list entry.
% Only textual answers are given here; drawings and individual observations
% that require the student's own work are not prescribed.

\section{Solutions}

% --- Models in Chemistry ---
\Solution{car-model}{%
  \textbf{a)} The models differ in material, scale and level of detail.

  \textbf{b)} All three represent a red car with a body and wheels.

  \textbf{c)} No. The best model depends on the question or purpose;
  each model highlights different features.
}

\Solution{wegener}{%
  \textbf{a)} Matching continental coastlines, identical fossils on
  widely separated continents and similar rock formations.

  \textbf{b)} Existence of former land bridges between continents.

  \textbf{c)} New observations may contradict a model or reveal its
  limitations, so scientists must revise or replace it.
}

\Solution{par-mod-fills}{%
  \textbf{a)} particles;
  \textbf{b)} properties (for example size or mass);
  \textbf{c)} empty space;
  \textbf{d)} attractive;
  \textbf{e)} faster (on average).
}

% --- Properties of Substances ---
\Solution{bar-chart}{%
  \textbf{a)} Sodium.

  \textbf{b)} Sodium -- lead -- aluminium -- copper -- manganese --
  nickel -- iron -- chromium.

  \textbf{c)} Copper: its melting point (about 1085~$^{\circ}$C) is
  closest to 1050~$^{\circ}$C among the metals shown.
}

\Solution{mp-table}{%
  \textbf{a)} Solids: sodium, cobalt. Liquids: ethanol, ethanoic acid.
  Gases: radon, nitrogen, propane.

  \textbf{b)} Radon, with a liquid temperature interval of
  $-62-(-71)=9~^{\circ}\mathrm{C}$.

  \textbf{c)} Radon and nitrogen.

  \textbf{d)} Nitrogen.

  \textbf{e)} Cobalt.
}

\Solution{sol-assign}{%
  A $\to$ 2; B $\to$ 5; C $\to$ 7; D $\to$ 1;
  E $\to$ 3; F $\to$ 6; G $\to$ 4.
}

\Solution{sol-graph}{%
  \textbf{a)} In the examples shown, many solid substances become more
  soluble as temperature rises, while gases such as NH$_3$ and SO$_2$
  become less soluble.

  \textbf{b)} The solubility of KNO$_3$ rises sharply with temperature;
  that solubility of NaCl changes only slightly.

  \textbf{c)} At room temperature (approximately 25~$^{\circ}$C), the
  graph indicates roughly 90~g of NaNO$_3$ per 100~g of water.
  This does not directly give the mass needed for 250~mL of
  \emph{saturated solution}: the graph specifies mass per mass of water,
  not per volume of the finished solution. The density of water would also be required.
}

\Solution{sol-sat}{%
  KCl is less soluble at room temperature than at 80~$^{\circ}$C.
  On cooling a solution that was saturated at 80~$^{\circ}$C, some
  dissolved KCl therefore crystallises out until the remaining
  solution reaches its new saturation concentration. At particle level,
  dissolved particles leave the solution and become arranged in a solid.
}

\Solution{dens-table}{%
  $\rho = 0.700~\mathrm{g/cm^3}$; $V=2.00~\mathrm{cm^3}$; $m=10.8~\mathrm{g}$.
}

\Solution{dens-cal}{%
  $\rho=1.2474~\mathrm{g/mL}=1247.4~\mathrm{kg/m^3}$.
}

\Solution{dens-vol}{%
  Gold: $\approx 518~\mathrm{cm^3}=0.518~\mathrm{L}$.
  Silver: $\approx 952~\mathrm{cm^3}=0.952~\mathrm{L}$.
}

\Solution{dens-unit}{%
  \textbf{a)} $19.3~\mathrm{g/cm^3}$;
  \textbf{b)} $19.3~\mathrm{kg/L}$;
  \textbf{c)} $19.3~\mathrm{mg/\mu L}$.
}

\Solution{dens-beam}{%
  The mercury side moves down. $V_{\mathrm{Hg}}=V_{\mathrm{Ag}} \to$
  $m_{\mathrm{Hg}}/m_{\mathrm{Ag}}=\rho_{\mathrm{Hg}}/
  \rho_{\mathrm{Ag}}=13.5/10.5\approx1.29$.
  Mercury therefore has approximately 29\,\% more mass than silver.
}

\Solution{dens-gly}{%
  $V\approx30.73~\mathrm{mL}$.
}

\Solution{dens-exp}{%
  \textbf{a)} $8.9~\mathrm{g/cm^3}$, copper.

  \textbf{b)} An attached air bubble also displaces water. The apparent
  volume is too large and the calculated metal density is too low.
}

\Solution{other-assign}{%
  A $\to$ 6; B $\to$ 5; C $\to$ 3;
  D $\to$ 7; E $\to$ 2; F $\to$ 9;
  G $\to$ 8; H $\to$ 1; I $\to$ 4.
}

% --- Further Applications of the Particle Model ---
\Solution{diff}{%
  Gas particles generally move faster and are much farther apart
  than liquid particles, allowing them to spread and mix more
  rapidly through random motion.
}

\Solution{mix-types}{%
  \textbf{Textual part only}:
  \textbf{1)} solid mixture (s/s);
  \textbf{2)} alloy (s/s);
  \textbf{3)} suspension (s/l);
  \textbf{4)} solution (s/l);
  \textbf{5)} solid aerosol (s/g);
  \textbf{6)} solid foam (g/s);
  \textbf{7)} emulsion (l/l);
  \textbf{8)} solution (l/l);
  \textbf{9)} liquid aerosol (l/g);
  \textbf{10)} foam (g/l);
  \textbf{11)} solution (g/l);
  \textbf{12)} gas mixture (g/g).
}

\Solution{mix-pure}{%
  \textbf{b)} graphite, \textbf{c)} table salt,
  \textbf{d)} diamond and \textbf{f)} dry ice.
}

\Solution{mix-hethomo}{%
  \textbf{a)} solid mixture (s/s).
  \textbf{b)} alloy (s/s).
  \textbf{c)} gas mixture (g/g).
}

\Solution{mix-assign}{%
  \textbf{1)} solid aerosol (s/g);
  \textbf{2)} liquid aerosol (l/g);
  \textbf{3)} alloy (s/s);
  \textbf{4)} solution (l/?);
  \textbf{5)} gas mixture (g/g);
  \textbf{6)} solution (l/?).
}

\Solution{mix-meri}{%
  Before: egg white and sugar $\to$ solution, egg white/sugar and air $\to$ foam; after: egg white/sugar and air $\to$ hard foam
}
